% Generated by review_v53/analyze_confirmation.py. Do not edit.
\begin{table}[!htb]
\centering\footnotesize
\setlength{\tabcolsep}{2pt}
\renewcommand{\arraystretch}{1.12}
\sisetup{group-digits=false}
\caption[Fester Bestätigungs-Durchsatz und gepaarte Intervalle]{Anhaltender Durchsatz der sieben festen Einsatzvergleiche. Full- und Split-FPS sind arithmetische Mittelwerte aus fünf Erfassungen mit frischem Prozess. Das Verhältnis ergibt sich aus dem Split-Mittelwert geteilt durch den Mittelwert der ungeteilten Ausführung. Nominelle 95\%-Fieller-Intervalle für gepaarte Werte verwenden die fünf ursprünglichen logischen Blöcke und vier Freiheitsgrade.}
\label{tab:confirmation:performance}
\begin{tabularx}{\linewidth}{@{}Y L{1.60cm}L{1.70cm}S[table-format=4.0]S[table-format=4.0]S[table-format=1.3]r@{}}
\toprule
Modell & Full & Split & {\shortstack{Full\\FPS}} & {\shortstack{Split\\FPS}} & {Verhältnis} & {95\%-Intervall} \\
\midrule
MobileNetV3-L & DX & DX / b135 & 1474 & 1603 & 1.087 & $[1.084,1.091]$ \\
RegNetX-1.6GF & DX & H10 / b123 & 542 & 1476 & 2.725 & $[2.698,2.751]$ \\
ResNet50 & DX & H10 / b095 & 660 & 1030 & 1.561 & $[1.547,1.575]$ \\
YOLO11l & DX & H10 / b120 & 63 & 53 & 0.838 & $[0.837,0.840]$ \\
YOLO26m & GPU (H8) & H10 / b043 & 82 & 110 & 1.352 & $[1.348,1.356]$ \\
YOLO26s & GPU (DX) & H8 / b021 & 141 & 189 & 1.337 & $[1.333,1.342]$ \\
YOLOv7 & DX & H8 / b066 & 83 & 113 & 1.360 & $[1.356,1.364]$ \\
\bottomrule
\end{tabularx}
\par\smallskip
\begin{minipage}{\linewidth}\scriptsize
DX steht für DEEPX DX-M1, H8 für Hailo-8 und H10 für Hailo-10H. GPU steht für TensorRT auf der Orin-GPU des in Klammern angegebenen Systems. Split-Einträge geben den Beschleuniger und den Schnitt an. Sieben der 35 logischen Durchsatzpaare erstrecken sich über mehrere Erfassungssitzungen. Die Intervalle basieren auf der Annahme unabhängiger, bivariat normalverteilter Blockbeobachtungen. Sie schätzen weder Sitzungsdrift noch bieten sie eine gleichzeitige Abdeckung. Tabelle~\ref{tab:confirmation:blocks} listet jeden Block auf.
\end{minipage}
\end{table}
